\documentclass[legalpaper,twocolumn]{article}
\usepackage[utf8]{inputenc}
\usepackage[T1]{fontenc}
\usepackage[spanish]{babel}
\usepackage{amsmath}
\usepackage[margin=0.5cm]{geometry}
\usepackage{titlesec}
\usepackage{enumitem}

\titleformat{\section}{\normalfont\large\bfseries}{\thesection}{1em}{}
\titleformat{\subsection}{\normalfont\normalsize\bfseries}{\thesubsection}{1em}{}

\setlength{\parindent}{0pt}
\setlength{\parskip}{3pt}

\title{\vspace{-1cm}Formulario - Sistemas de Control II}
\author{Emerson Warhman}
\date{}

\begin{document}
\maketitle
\vspace{-1.5cm}

\section*{Sistemas de Segundo Orden}
Función de transferencia (lazo cerrado):
\[ T(s) = \frac{\omega_n^2}{s^2 + 2\zeta \omega_n s + \omega_n^2} \]
Polos:
\[ s = -\zeta \omega_n \pm \omega_n \sqrt{\zeta^2 - 1} \]
\textbf{Transitorio (Subamortiguado $0 < \zeta < 1$):}
\begin{itemize}[leftmargin=*]
    \item Tiempo de alza: \( t_r = \frac{\pi - \tan^{-1}(\frac{\sqrt{1-\zeta^2}}{\zeta})}{\omega_n} \)
    \item Tiempo pico: \( t_p = \frac{\pi}{\omega_n \sqrt{1-\zeta^2}} \)
    \item Sobrepico: \( M_p = e^{-\frac{\zeta \pi}{\sqrt{1-\zeta^2}}} \)
    \item Tiempo de establecimiento (2\%): \( t_s = \frac{4}{\zeta \omega_n} \)
\end{itemize}

\section*{Estabilidad}
\subsection*{Criterio de Jury (Discreto)}
Polinomio: \( P(z) = a_n z^n + \dots + a_1 z + a_0 = 0 \)
\begin{enumerate}[leftmargin=*]
    \item Cond. necesaria: \( |a_0| < a_n \) (o $|a_0| < 1$ si $a_n=1$).
    \item Tabla de Jury:
    \[
    \begin{array}{c|ccccc}
        F_1 & a_n & a_{n-1} & \dots & a_0 \\
        F_2 & a_0 & a_1 & \dots & a_n \\
        F_3 & b_{n-1} & b_{n-2} & \dots & \\
        \dots & & & & 
    \end{array}
    \]
    donde \( b_k = \det \begin{pmatrix} a_n & a_k \\ a_0 & a_{n-k} \end{pmatrix} \) (verificar def. exacta).
    \item Cond. suficiente: Todos los elementos de la 1ra columna $>0$.
\end{enumerate}

\subsection*{Criterio de Routh (Continuo)}
Polinomio: \( s^n + a_{n-1}s^{n-1} + \dots + a_0 = 0 \)
Elemento $b_1$:
\[ b_1 = -\frac{1}{a_{k-1,1}} \det \begin{pmatrix} a_{k-2,1} & a_{k-2,2} \\ a_{k-1,1} & a_{k-1,2} \end{pmatrix} \]
Condición: No cambios de signo en 1ra columna.

\section*{Lugar Geométrico de las Raíces (LGR)}
\textbf{Pasos para dibujar:}
\begin{enumerate}[leftmargin=*,nosep]
    \item Polos y ceros de G(s)H(s).
    \item Eje Real: Regla de paridad (polos+ceros a la derecha es impar).
    \item Asíntotas:
    \[ \theta_a = \frac{180(2k+1)}{N_p - N_z}, \quad \sigma_a = \frac{\sum P - \sum Z}{N_p - N_z} \]
    \item Ruptura/Ingreso: Resolver $\frac{dK}{ds} = 0$. ($K = -1/GH(s)$).
    \item Ángulos de Salida/Llegada:
    \[ \theta_{sal} = 180 - \sum \angle P_{otros} + \sum \angle Z \]
    \[ \theta_{lleg} = 180 - \sum \angle Z_{otros} + \sum \angle P \]
    \item Intersección Eje Img: Routh o $s=j\omega$.
    \item Trazar uniendo polos a ceros/infinito.
\end{enumerate}

\section*{Compensadores}
\textbf{Adelanto (Lead)} ($\alpha < 1$):
\[ G_c(s) = K_c \frac{s + 1/T}{s + 1/(\alpha T)} \]
Diseño por LGR:
\begin{enumerate}[leftmargin=*,nosep]
    \item Definir punto deseado $s_d$ (según $\zeta, \omega_n$).
    \item Calcular deficiencia angular: $\phi = 180 - \sum \angle \text{vec}(s_d)$.
    \item Ubicar polo $p$ y cero $z$ tal que aporten $\phi$.
    \item Fórmula bisectriz o usar $\alpha = \frac{1-\sin \phi_m}{1+\sin \phi_m}$ y ajustar $T$.
    \item Calcular $K_c$ por condición de magnitud en $s_d$.
\end{enumerate}

\textbf{Atraso (Lag)} ($\beta > 1$):
\[ G_c(s) = K_c \frac{s + 1/T}{s + 1/(\beta T)} \]
Diseño por LGR:
\begin{enumerate}[leftmargin=*,nosep]
    \item Ajustar ganancia $K$ para cumplir requisitos transitorios (polos dominantes s/compensar).
    \item Calcular constante error actual y la requerida $\to$ hallar $\beta$.
    \item Ubicar cero $z$ cerca del origen (ej. $0.1 Re(s_d)$) y polo $p = z/\beta$.
    \item Ajustar $K_c \approx K$.
\end{enumerate}

\textbf{Controlador PID}:
\[ G_c(s) = K_p + \frac{K_i}{s} + K_d s \]

\section*{Sistemas Muestreados}
Transformación: \( z = e^{sT} \), \( s = \frac{1}{T} \ln z \) \\
Mapeo de polos:
\[ |z| = e^{-\zeta \omega_n T}, \quad \angle z = \omega_d T \]
donde $\omega_d = \omega_n \sqrt{1-\zeta^2}$.

\section*{Espacio de Estados}
\textbf{Variables de Fase (Forma Canónica)}
Ecuación dif: \( y^{(n)} + \dots + a_0 y = b_m u^{(m)} + \dots + b_0 u \)
Matriz A (común):
\[
A = \begin{bmatrix}
0 & 1 & 0 & \dots \\
\vdots & & \ddots & \\
-a_0 & -a_1 & \dots & -a_{n-1}
\end{bmatrix}
\]
\textbf{Caso 1: Ceros en numerador (Modificar C)}
\[ G(s) = \frac{K(c_m s^{m-1} + \dots + c_1)}{s^n + \dots} \]
\[ B = [0, \dots, 0, K]^T \]
\[ C = [c_1, c_2, \dots, c_m, 0, \dots] \]
\textbf{Caso 2: Ceros en numerador (Modificar B)}
Usar unicidad de G(s): $G(s) = C(sI-A)^{-1}B$.
\[ C = [1, 0, \dots, 0] \]
B se halla igualando coeficientes:
\[ \text{Num}(C(sI-A)^{-1}B) = \text{Num}(G(s)) \]

\subsection*{Matrices Generales}
Controlabilidad: $\mathcal{C} = [B, AB, \dots, A^{n-1}B]$ \\
Observabilidad: $\mathcal{O} = [C^T, (CA)^T, \dots, (CA^{n-1})^T]^T$

\subsection*{Forma Modal (Diagonal)}
Si $A$ tiene autovalores distintos $\lambda_i$:
\[ \Lambda = P^{-1} A P = \text{diag}(\lambda_1, \dots, \lambda_n) \]
donde $P$ contiene los autovectores.
\[ \dot{z} = \Lambda z + P^{-1} B u , \quad y = C P z \]
Expansión fracciones parciales: $G(s) = \sum \frac{d_i}{s-\lambda_i}$
\end{document}
